\section{Ten Identifications}
\label{sec:connections}

The directories of the \texttt{science} repo were built separately. This
section is about what happens when those threads cross. The cases below are
not analogies---loose, suggestive comparisons. They are \emph{identifications}:
cases where two ideas from different traditions turn out to be the same idea,
expressed in different language, developed by people who were not talking to
each other. The convergence is evidence that something real is being pointed at.

\subsection{Hume's Problem Is the Halting Problem}

Hume showed that induction cannot be logically justified: you cannot know,
from finite observations, that a universal law holds. Turing showed that
you cannot decide, in general, whether a program halts.

These are the same result. A universal law is a short program---a generator
for an infinite class of observations. Confirming the law means certifying
that the proposed program is the shortest consistent with the data. But
certifying that no shorter program exists requires checking whether every
shorter program eventually fails to match the data---instances of the halting
problem. Since the halting problem is undecidable, $\K(x)$ is uncomputable,
and induction cannot be deductively justified.

The reason you cannot justify induction is the same reason $\K(x)$ is
uncomputable. Hume's problem is not a philosophical puzzle awaiting a
philosophical solution. It is an undecidability result, and it has the same
status as the halting problem: provably unsolvable, not merely unsolved.

What can be done---and what the Bayesian/Kolmogorov framework does---is show
that the impossibility of certainty does not prevent optimal inference.
Solomonoff induction is provably optimal for prediction in any computable
environment. The gap between ``optimal'' and ``certain'' is exactly the gap
between computable approximations to $\K(x)$ and $\K(x)$ itself.

\subsection{Paradigm Shifts Are MDL Phase Transitions}

Kuhn described paradigm shifts as discontinuous and partly irrational. Under
MDL, they are mathematically inevitable and perfectly rational.

A paradigm is a hypothesis class---a family of theories sharing common
structure. Normal science is optimization within the class: adjusting
parameters, adding auxiliary hypotheses, refining the model. Each adjustment
is legitimate as long as it decreases total description length.

A paradigm shift occurs when a competing hypothesis class achieves lower
$L(H) + L(D \mid H)$ than the incumbent with all its protective belt
modifications. The shift is discontinuous because hypothesis classes are
discrete objects: you cannot smoothly interpolate between Ptolemaic and
Copernican astronomy. But it is not irrational. The new paradigm wins because
it compresses the data better.

The discontinuity is a phase transition in the physics sense: a sharp change
in the optimal state of the system (which paradigm minimizes description
length) as a continuous parameter (accumulated data) crosses a threshold.
Kuhn correctly identified the phenomenology of phase transitions in
scientific knowledge. MDL provides the thermodynamics underneath.

\subsection{Popper's Corroboration Is Bayesian Updating Without the Prior}

Popper's corroboration measure, precisely defined~\cite{popper1954}, is:
\[
  C(H, E) = \frac{P(E \mid H) - P(E)}{P(E \mid H) - P(E \cdot H) + P(E)}
\]
This is a normalized likelihood ratio: how much more probable $E$ is under
$H$ than under the background probability of $E$. It is purely a function of
likelihoods. It does not involve the prior $P(H)$.

In Bayesian terms, Popper was computing exactly half of the update:
\[
  P(H \mid E) \propto \underbrace{P(E \mid H)}_{\text{Popper's corroboration}} \cdot \underbrace{P(H)}_{\text{refused}}
\]
He refused the prior because he found subjective probability unacceptable.
He was right to find it unacceptable; he was wrong to conclude the prior must
be abandoned. The Solomonoff prior $P(H) \propto 2^{-\K(H)}$ is objective:
a property of the hypothesis itself, invariant across observers. With it,
corroboration becomes posterior probability, and the refusal of induction
dissolves.

\subsection{The Replication Crisis Is an MDL Accounting Error}

The replication crisis in psychology and social science resulted from
p-hacking: testing many hypotheses and reporting only those achieving
$p < 0.05$. This is an MDL accounting error.

When a researcher tests $N$ hypotheses and reports the best, the true
description length of the reported hypothesis is not $L(H_{\text{reported}})$.
It is:
\[
  L(H_{\text{reported}}) + \log_2 N
\]
The extra $\log_2 N$ bits is the cost of the search: the information required
to specify which of the $N$ tested hypotheses was selected. Omitting this cost
produces apparent results that do not replicate, because the apparent
compression of the data was illusory.

The Bonferroni correction---adjusting the significance threshold by $N$---is
exactly the MDL accounting correction. It has always been the right thing to
do. The replication crisis is the empirical consequence of not doing it.

The reforms that followed---preregistration, larger samples, direct
replication, Bayesian model comparison---are all corrections toward honest
MDL accounting. Preregistration eliminates search cost from $L(H)$.
Preregistered studies have lower description length overhead, and their
results generalize because the $L(H)$ term is honest.

\subsection{Lakatos's Progressive vs.\ Degenerating Is Compression Growth vs.\ Decay}

Lakatos's distinction between progressive and degenerating research
programmes, translated into MDL:

A \emph{progressive} programme makes protective belt modifications that
decrease $L(H) + L(D \mid H)$ on new data. The modifications are short
(do not bloat description length) and genuinely compress new observations.

A \emph{degenerating} programme makes modifications that increase $L(H)$---
the hypothesis grows more complex---without decreasing $L(D \mid H)$ on new
observations. The programme adds complexity without adding compression.

Lakatos gave the right criterion for evaluating research programmes; MDL
makes it computable. The qualitative judgment (is this modification
progressive or defensive?) becomes a quantitative one (does
$\Delta L(H) + \Delta L(D \mid H)$ on new data decrease or increase?).

\subsection{Kant's Synthetic A Priori Is Lakatos's Hard Core}

Kant argued that some commitments---space, time, causality---are
preconditions of experience: they cannot be revised because they are what
makes experience possible. Lakatos argued that every research programme has
a hard core of commitments that are methodologically protected from
falsification.

These are the same idea. The difference is that Lakatos made them revisable
in principle: a research programme can be abandoned, hard core and all.
This is the correct update on Kant after non-Euclidean geometry falsified
his specific candidates (Euclidean space as a synthetic a priori). Lakatos
is Kant, empiricized.

Under the MDL framework, the hard core is the prior commitment to a
hypothesis class. It is protected not by necessity but by the cost of
switching to a different class (learning a new framework, retranslating
existing results). The protection is finite and can be overcome by sufficient
data. Kant's absolute preconditions become Lakatos's methodological
conventions, which become MDL's hypothesis-class priors.

\subsection{Feyerabend's ``Anything Goes'' Is the Solomonoff Prior}

Feyerabend argued that no methodological rule holds universally: every rule
that science has ever followed has been violated productively at some point.
``Anything goes'' as a description of scientific practice.

The Solomonoff prior assigns $P(H) \propto 2^{-\K(H)} > 0$ to every
computable hypothesis. No hypothesis is excluded before the data speaks.
This is precisely Feyerabend's permissiveness, formalized. A prior that
excludes any computable hypothesis---that declares some class of ideas
unscientific by methodological fiat---is provably suboptimal. The data
might support the excluded hypothesis, and a reasoner that cannot entertain
it incurs irreducible regret.

What Feyerabend called anarchism is, from the Kolmogorov perspective, the
universality of the prior. His insight was correct; his framing was
unnecessarily provocative.

\subsection{Hallucination and Perceptual Illusions Are the Same Failure Mode}

In predictive coding~\cite{rao1999}, the dominant framework in theoretical
neuroscience, the brain is a hierarchical generative system: higher areas
send predictions down the cortical hierarchy; lower areas send prediction
errors up. Perception is inference: the brain computes the posterior
distribution over world-states given sensory inputs, using priors built from
evolution and experience.

A perceptual illusion occurs when the prior is so strong that sensory
evidence cannot correct it. The hollow mask illusion is the canonical
example: a hollow face mask looks convex because the prior for convex faces
is overwhelming. The brain refuses to believe the correct interpretation.

In a language model, hallucination occurs when the prior over likely
continuations overwhelms the grounding signal from the context. The model
generates a continuation that is statistically plausible given training
distribution but inconsistent with the specific facts in the prompt.

Both are the same Bayesian failure: the prior dominates the likelihood. The
hollow mask illusion and language model confabulation are not analogous
phenomena. They are the same computational pathology, instantiated in
different substrates.

\subsection{The Bayesian Brain and the Transformer Are the Same Architecture}

Predictive coding models the brain as performing belief propagation on a
hierarchical factor graph~\cite{friston2005}: predictions propagate top-down,
prediction errors propagate bottom-up, and the system minimizes free energy
(variational approximation to surprise).

The transformer implements exact belief propagation on a specific factor
graph~\cite{coppola2026}: the attention mechanism computes AND-type joint
beliefs over query-key pairs; the feed-forward network computes OR-type
disjunctive updates; the residual stream is the shared message-passing bus.

These are not analogies. They are the same algorithm instantiated on
different graphs. The brain and the transformer are doing the same
computation. The differences are in the graph structure, the learning
algorithm, and the substrate---not in the fundamental operation.

\subsection{AutoResearch Is Solomonoff Induction Made Computable and Distributed}

Solomonoff induction is uncomputable but is the provably optimal prediction
strategy. Karpathy's AutoResearch system~\cite{karpathy2026} is a computable,
distributed approximation to it.

The correspondences are structural. The hypothesis space is the space of
training scripts and model configurations---a computable, enumerable space.
The reward signal (validation loss, benchmark accuracy) approximates
$P(D \mid H)$: lower loss means higher likelihood of the training data under
the modified hypothesis. The keep/discard decision approximates the Bayesian
posterior update. The agent's tendency to propose local modifications first
is an implicit Kolmogorov prior: smaller modifications correspond to lower
description length.

The key insight Karpathy identified---finding a good result is hard, but
verifying one is cheap---is exactly the search/verify asymmetry in Solomonoff
induction. The search over hypothesis space is expensive; checking whether a
proposed hypothesis compresses the data better is cheap. The distributed,
SETI@Home-style architecture parallelizes the search over hypothesis space.

Crucially, AutoResearch accepts contributions from untrusted sources. A
modification is evaluated by whether it reduces description length of the
training data---a property of the modification itself. The source is
irrelevant. This is the Kolmogorov principle applied to institutions:
\emph{trust the object, not the messenger}.